% Part: normal-modal-logic
% Chapter: filtrations
% Section: more-filtrations

\documentclass[../../../include/open-logic-section]{subfiles}

\begin{document}

\olfileid{nml}{fil}{acc}

\olsection{వడపోతలు మరియు ప్రాప్యత ధర్మాలు}

ముందు చెప్పినట్లు, సార్వత్రిక, సీరియల్,
స్వావర్తన నమూనాల వడపోతలు కూడా వరుసగా సార్వత్రిక,
సీరియల్, స్వావర్తన నమూనాలే. కానీ సౌష్ఠవ లేదా
సంక్రామక నమూనాకు చేసే ప్రతి వడపోత వరుసగా సౌష్ఠవ
లేదా సంక్రామక నమూనా కావాల్సిన అవసరం లేదు.
అయితే కొన్ని సందర్భాల్లో ఆ ధర్మాలు నిలిచేలా
వడపోతలను నిర్వచించవచ్చు. అందుకోసం అత్యంత స్థూల
వడపోత నిర్వచనం మాదిరిగానే చేసి,~$R^*$ నిర్వచనానికి
అదనపు షరతులను చేర్చుతాం. $\Gamma$ \tetoken{ఉపసూత్రాల}{!!{formula}s}
పరంగా సంవృతమై ఉండనివ్వండి. నమూనా
$\mModel{M} =\tuple{W, R, V}$లోని లోకాలు $u$, $v$
మధ్య \olref{tab:Cn-filtrations}లో ఉన్న
సంబంధాలు~$C_i(u,v)$ను పరిశీలిద్దాం.
ఈ షరతుల సంయోజనాల ఆధారంగా $R^*[u][v]$ను
నిర్వచించవచ్చు. ఉదాహరణకు, $C_1(u,v)$
నెరవేరితే, అప్పుడు మాత్రమే $R^*[u][v]$ అని
నిర్దేశిస్తే, ఖచ్చితంగా అత్యంత స్థూల వడపోతే
వస్తుంది. $C_1(u,v)$, $C_2(u, v)$ రెండూ
నెరవేరితే, అప్పుడు మాత్రమే $R^*[u][v]$ అని
నిర్దేశిస్తే, వేరొక వడపోత వస్తుంది. కేవలం
$C_1(u,v)$ను నెరవేర్చే జతలకంటే రెండింటినీ
$C_1(u,v)$, $C_2(u,v)$ నెరవేర్చే లోకాల జతలు
తక్కువ కనుక ఇది అత్యంత స్థూల వడపోతకంటే
``సూక్ష్మమైనది''.

\begin{table}[ht]
  \centering
  \begin{tabular}{|ll|}
    \hline
    \multirow{2}{*}{$C_1(u,v)$:}
    \iftag{prvBox}{&
      $\Box!A \in \Gamma$, $\mSat{M}{\Box!A}[u]$ అయితే
      $\mSat{M}{!A}[v]$;\iftag{prvDiamond}{ మరియు}{}\\}{}
    \iftag{prvDiamond}{&
      $\Diamond!A \in \Gamma$, $\mSat{M}{!A}[v]$ అయితే
      $\mSat{M}{\Diamond!A}[u]$; \\}{}
    \hline
    \multirow{2}{*}{$C_2(u,v)$:}
    \iftag{prvBox}{&
      $\Box!A \in \Gamma$, $\mSat{M}{\Box!A}[v]$ అయితే
      $\mSat{M}{!A}[u]$;\iftag{prvDiamond}{ మరియు}{}\\}{}
    \iftag{prvDiamond}{&
      $\Diamond!A \in \Gamma$, $\mSat{M}{!A}[u]$ అయితే
      $\mSat{M}{\Diamond!A}[v]$; \\}{}
    \hline
    \multirow{2}{*}{$C_3(u,v)$:}
    \iftag{prvBox}{&
      $\Box!A \in \Gamma$, $\mSat{M}{\Box!A}[u]$ అయితే
      $\mSat{M}{\Box!A}[v]$;\iftag{prvDiamond}{ మరియు}{}\\}{}
    \iftag{prvDiamond}{&
      $\Diamond!A \in \Gamma$, $\mSat{M}{\Diamond!A}[v]$ అయితే
      $\mSat{M}{\Diamond!A}[u]$; \\}{}
    \hline
    \multirow{2}{*}{$C_4(u,v)$:}
    \iftag{prvBox}{&
      $\Box!A \in \Gamma$, $\mSat{M}{\Box!A}[v]$ అయితే
      $\mSat{M}{\Box!A}[u]$;\iftag{prvDiamond}{ మరియు}{}\\}{}
    \iftag{prvDiamond}{&
      $\Diamond!A \in \Gamma$, $\mSat{M}{\Diamond!A}[u]$ అయితే
      $\mSat{M}{\Diamond!A}[v]$; \\}{}
    \hline
    \end{tabular}
  \caption{వడపోతలను నిర్వచించడానికి లోకాలపై షరతులు.}
  \ollabel{tab:Cn-filtrations}
\end{table}

\begin{thm}\ollabel{thm:more-filtrations}
  $\mModel{M} =\tuple{W,R,V}$ ఒక నమూనా,
  $\Gamma$ \tetoken{ఉపసూత్రాల}{!!{formula}s} పరంగా
  సంవృతమై ఉండనివ్వండి. $W^*$, $V^*$లు
  \olref[fil]{defn:filtration}లోలాగే
  నిర్వచించబడ్డాయని అనుకుందాం. అప్పుడు:
  \begin{enumerate}
  \item $C_1(u, v) \land C_2(u,v)$ అయితే, అప్పుడు
    మాత్రమే $R^*[u][v]$ అనుకుందాం. అప్పుడు $R^*$
    సౌష్ఠవం; $\mModel{M}$ సౌష్ఠవమైతే
    $\mModel{M^*} = \tuple{W^*,R^*,V^*}$ ఒక వడపోత.
  \item $C_1(u, v) \land C_3(u,v)$ అయితే, అప్పుడు
    మాత్రమే $R^*[u][v]$ అనుకుందాం. అప్పుడు $R^*$
    సంక్రామకం; $\mModel{M}$ సంక్రామకమైతే
    $\mModel{M^*}=\tuple{W^*,R^*,V^*}$ ఒక వడపోత.
  \item $C_1(u, v) \land C_2(u,v)
    \land C_3(u,v) \land C_4(u,v)$ అయితే,
    అప్పుడు మాత్రమే $R^*[u][v]$ అనుకుందాం.
    అప్పుడు $R^*$ సౌష్ఠవమూ సంక్రామకమూ;
    $\mModel{M}$ సౌష్ఠవమూ సంక్రామకమూ అయితే
    $\mModel{M^*}=\tuple{W^*,R^*,V^*}$ ఒక వడపోత.
  \item $C_1(u,
    v) \land C_3(u,v) \land C_4(u,v)$ అయితే,
    అప్పుడు మాత్రమే $R^*[u][v]$ అనేలా $R^*$ను
    నిర్వచించామనుకుందాం. అప్పుడు $R^*$ సంక్రామకమూ
    యూక్లిడియన్‌దీ; $\mModel{M}$ సంక్రామకమూ
    యూక్లిడియన్‌దీ అయితే
    $\mModel{M^*}=\tuple{W^*,R^*,V^*}$ ఒక వడపోత.
  \end{enumerate}
\end{thm}

\begin{proof}
  \begin{enumerate}
    \item $C_1(u,v)
      \Leftrightarrow C_2(v,u)$,
      $C_2(u,v) \Leftrightarrow C_1(v,u)$ కనుక
      $R^*$ సౌష్ఠవమని తక్షణమే తెలుస్తుంది.
      కాబట్టి $\mModel{M}$ సౌష్ఠవమైతే
      $\mModel{M^*}$,~$\Gamma$ ద్వారా వడపోత
      అని చూపడమే మిగిలింది. $C_1(u,v)$ వల్ల
      \iftag{prvBox}{%
        \iftag{prvDiamond}{%
          \olref[fil]{defn:filtration-R2} మరియు
          \olref[fil]{defn:filtration-R3} అనే
          \olref[fil]{defn:filtration} షరతులు}{%
          \olref[fil]{defn:filtration}\olref[fil]{defn:filtration-R2} షరతు}}{%
        \olref[fil]{defn:filtration}\olref[fil]{defn:filtration-R3} షరతు}
      నెరవేరినట్లే. కనుక ఇక
      \olref[fil]{defn:filtration}\olref[fil]{defn:filtration-R1},
      అంటే $Ruv$ నుంచి $R^*[u][v]$ వస్తుందని
      నిర్ధారించాలి.

      $Ruv$ అనుకుందాం. $R^*[u][v]$ కోసం
      $C_1(u,v)$, $C_2(u,v)$లను స్థాపించాలి.
      $C_1$ కోసం: \iftag{prvBox}{$\Box!A
        \in\Gamma$, $\mSat{M}{\Box!A}[u]$ అయితే
        $Ruv$ కనుక $\mSat{M}{!A}[v]$ కూడా.
        \iftag{prvDiamond}{అలాగే, }{}}{}\iftag{prvDiamond}{$\Diamond!A \in \Gamma$,
        $\mSat{M}{!A}[v]$ అయితే $Ruv$ కనుక
        $\mSat{M}{\Diamond!A}[u]$.}{}
      $C_2$ కోసం: \iftag{prvBox}{$\Box!A \in \Gamma$,
        $\mSat{M}{\Box!A}[v]$ అయితే, సౌష్ఠవం వల్ల
        $Ruv$ నుంచి $Rvu$; అందువల్ల
        $\mSat{M}{!A}[u]$.\iftag{prvDiamond}{అలాగే, }{}}{}\iftag{prvDiamond}{$\Diamond!A \in\Gamma$,
        $\mSat{M}{!A}[u]$ అయితే, సౌష్ఠవం వల్ల
        $Rvu$ కనుక $\mSat{M}{\Diamond!A}[v]$.}{}
    \item వ్యాయామం.
    \item వ్యాయామం.
    \item వ్యాయామం.
  \end{enumerate}
\end{proof}

\begin{prob}
  \olref[nml][fil][acc]{thm:more-filtrations} నిరూపణను
  పూర్తి చేయండి.
\end{prob}

\end{document}
